\documentclass[11pt,a4paper]{article}
\usepackage[T1]{fontenc}
\usepackage[utf8]{inputenc}
\usepackage{lmodern}
\usepackage{amsmath,amssymb}
\usepackage[margin=25mm]{geometry}
\usepackage[hidelinks]{hyperref}
\hypersetup{pdftitle={Score-constrained ensembles and population transport},pdfauthor={navstokgap research working notes}}
\newcommand{\tightlist}{\setlength{\itemsep}{0pt}\setlength{\parskip}{0pt}}
\setlength{\emergencystretch}{3em}
\setcounter{secnumdepth}{0}
\begin{document}
\section{Score-constrained ensembles and population
transport}\label{score-constrained-ensembles-and-population-transport}

Restricting the ordinary action of two classical phase sheets to the
shared-score preparations gives the canonical Fisher field action
exactly. The restriction is not preserved by free Newtonian transport. A
separate population-transfer construction identifies its full
transport-conjugate action; bridge phase alone omits reservoir costs.
The resulting two-coordinate model satisfies continuity, but its shape
family fails the full field equation, even after adding Gaussian width.

These are variational identities and a stated closure obstruction,
developed by GPT-6.1 Sol and GPT-6 Astra on 2026-10-01. They do not
derive a physical constraint mechanism, quantum recording, or positive
action necessity. The supplied parameter is \(\kappa\ge0\); the field
obstruction below concerns \(\kappa>0\). Proof status: written,
internally checked.

\subsection{1. What the classical sheet action
supplies}\label{what-the-classical-sheet-action-supplies}

Work on a connected position interval or the line. Require smooth
positive densities, finite displayed energies, and vanishing boundary
terms for the variations used below. A classical phase sheet has
position density \(r\) and momentum \(W_q\). Its canonical one-form and
free Hamiltonian are \(\int W\,\delta r\,dq\) and
\(\int rW_q^2/(2m)dq\), respectively, with \(m>0\).

This one-form follows from the ordinary particle form. If material
coordinates \(a\) carry fixed probability mass, \(q=X(a)\),
\(r\,dq=da\), and \(v=\delta X\circ X^{-1}\) is a virtual displacement,
then \(\delta r=-\partial_q(rv)\). Integration by parts gives
\[\int W_q(X(a))\delta X(a)da
 =\int rW_qv\,dq=\int W\,\delta r\,dq.\tag{1}\] Fixed total mass makes a
spatially constant shift of \(W\) a gauge. This is a smooth variational
identity; no infinite-dimensional well-posedness theorem is being
assumed.

For two sheets of probability mass \(1/2\), introduce the constraint
\[r_\sigma=\rho/2,\qquad
W_\sigma=S+\kappa\sigma\log(\rho/\rho_{\rm ref}),\qquad
\sigma\in\{+1,-1\}.\tag{2}\] Here \(\rho\) is normalized, \(S\) has
action units, and the positive constant \(\rho_{\rm ref}\) only makes
the logarithm dimensionless. Changing it adds sheetwise constants and
changes no momentum. Formula (2) is the field version of the supplied
positive classical preparation in
\href{sed-closure-under-recording.md\#822-a-shared-score-sign-realizes-copying-but-fails-two-completion-tests}{the
recording note, §8.22}.

\textbf{Proposition 1 (exact constrained action).} Pulling the two-sheet
one-form and energy back to (2) gives \[\begin{aligned}
\Theta&=\sum_\sigma\int W_\sigma\delta r_\sigma\,dq
       =\int S\delta\rho\,dq,\\
\mathcal H&=\sum_\sigma\int\frac{r_\sigma W_{\sigma,q}^2}{2m}dq
 =\int\frac{\rho S_q^2}{2m}dq+\frac{\kappa^2}{2m}I(\rho),\\
I(\rho)&=\int\frac{\rho_q^2}{\rho}dq.
\end{aligned}\tag{3}\] The first identity cancels the two score signs;
the second expands the two squares. Thus variation of the
\emph{restricted} action \(\int[\int S\rho_t\,dq-\mathcal H]dt\) yields
\[\rho_t=-\partial_q(\rho S_q/m),\qquad
S_t=-\frac{S_q^2}{2m}
 +\frac{2\kappa^2}{m}\frac{(\sqrt\rho)_{qq}}{\sqrt\rho}.\tag{4}\] No
further canonical-field bracket is needed after adopting the constraint
and restricted variation. The Fisher modification is established prior
art {[}@HallReginatto2002; reading: passages pp.~3--8, especially the
classical action (4) and modified action/equations (17)--(18);
\href{../docs/batches/B09/Hall_Reginatto_2002.md}{source companion}{]}.
The calculation here identifies which correlated sheet restriction
produces that action and tests its invariance.

\textbf{Why this is not an invariant reduction.} Unrestricted classical
sheet continuity gives, on (2), \[\partial_t(r_+-r_-)
 =-\frac\kappa m\partial_q^2\rho.\tag{5}\] For \(\kappa>0\), no strictly
positive normalized density on the whole line has \(\rho_{qq}\equiv0\):
an affine function cannot have those properties. Consequently every such
constrained preparation departs from equal sheet densities immediately
somewhere. Restricted variation and ordinary free sheet transport are
different dynamics. Additional constraint reactions or state changes are
required, as the exact Gaussian switching test in
\href{sed-closure-under-recording.md\#823-a-stationary-switching-process-repairs-the-gaussian-marginal}{§8.23}
also demonstrates. In particular (4)'s current uses the mean momentum
\(S_q\); it does not give both physical sheet velocities
\(W_{\sigma,q}/m\) without further dynamics.

On connected positive support the labeled sheet momenta determine
\(S_q\) and hence \(S\) up to a common constant. At exact vacuum,
separate constants on disconnected components become invisible to the
ordinary physical phase-space measure. This construction cannot carry
them into a vacuum limit merely by declaring them canonical data.

\subsection{2. Population transport includes the
reservoirs}\label{population-transport-includes-the-reservoirs}

Let \(h_L,h_R\) be nonnegative smooth normalized bumps with compact,
ordered, disjoint supports. Between them fix a bridge \(B=[q_L,q_R]\).
Choose a smooth positive normalized density \(\rho_0\) with finite
Fisher information. Put \[\begin{aligned}
h&=h_R-h_L,&\rho_w&=\rho_0+wh,\\
G(q)&=\int_{-\infty}^{q}(h_L-h_R)(y)dy,&
M(w)&=\int\frac{G(q)^2}{\rho_w(q)}dq,\\
R_B&=\int_B\frac{dq}{\rho_0(q)}.&&
\end{aligned}\tag{6}\] Here \(w\) is probability mass transferred from
left to right, \(0\le G\le1\), \(G=1\) on \(B\), and \(G\) vanishes
outside the reservoirs and their connector. If \(K\) is their compact
enclosing interval, \(c=\min_K\rho_0>0\) and \(H=\|h\|_\infty>0\), then
\(|w|<c/(2H)\) ensures \(\rho_w\ge c/2\) on \(K\). Thus all terms in (6)
are finite, with \(M(w)\le2|K|/c\), and smooth in this \(w\) range.

Vanishing exterior flux and continuity uniquely require
\(j(q)=\dot w\,G(q)\), because \(\partial_qj=-\partial_t\rho_w\). The
corresponding total transport kinetic energy is \(mM(w)\dot w^2/2\).
Both \(M\) and \(R_B\) have length-squared units. Smooth reservoir
transport contributes positively, so \(M(w)>R_B\).

\textbf{Proposition 2 (continuity-compatible canonical family).} Define
an action field, up to a common constant, by
\[S_q(q;w,P)=\frac{P G(q)}{M(w)\rho_w(q)}.\tag{7}\] The variable \(P\)
has action units; it is not a point-particle momentum. Then
\[\begin{aligned}
\Theta&=\int S\delta\rho_w\,dq=P\,dw,\\
H(w,P)&=\frac{P^2}{2mM(w)}+\frac{\kappa^2}{2m}I(\rho_w),\\
\dot w&=\frac{P}{mM(w)},&
\dot P&=\frac{P^2M'(w)}{2mM(w)^2}
       -\frac{\kappa^2}{2m}\frac{dI(\rho_w)}{dw}.
\end{aligned}\tag{8}\] These local reduced Hamilton equations preserve
\(H\) while \(w\) remains in the positive-density range and satisfy the
ambient continuity equation exactly.

\emph{Proof.} Since \(h=-G_q\), integration by parts yields
\(\int S h\,dq=\int S_qG\,dq=P\). The phase energy in (3) is
\(P^2/(2mM)\), proving the Hamiltonian and its equations. Their first
equation gives \(\rho_wS_q/m=\dot wG\), the required current.
\(\square\)

In this family the actual bridge action difference is
\[\Phi_B=S(q_R)-S(q_L)=\frac{R_B}{M(w)}P.\tag{9}\] Thus \(\Phi_B\)
generally differs from the population-conjugate action
\(P=\int S h_Rdq-\int S h_Ldq\). The latter is a reservoir-weighted
phase-action difference. Replacing it by bridge endpoint phases omits
the reservoir motion. In coordinates \((w,\Phi_B)\) the canonical
two-form has coefficient \(M(w)/R_B\), rather than one.

For a fixed bridge density the exact dual identities are
\[E_{{\rm transport},B}=\frac{mJ^2R_B}{2},\qquad
\inf E_{{\rm phase},B}=\frac{\Phi_B^2}{2mR_B},\qquad
J=\frac{\Phi_B}{mR_B}.\tag{10}\] The phase infimum is
\href{sed-closure-under-recording.md\#821-the-exact-energy-of-a-phase-bridge-and-its-cut-law}{the
bridge lemma, (131)}, written with \(S\) instead of \(\hbar\theta\); it
is attained in the stated weighted \(H^1\) class. The last identity
makes the two energies equal. A low-density connector makes a fixed
phase difference cheap and a fixed population flux expensive. It selects
no absolute action scale.

The costs add over arbitrary spatial cuts, because both integrals of
\(1/\rho\) and of \(G^2/\rho\) add. This is an exact spatial transport
identity, not Newton's temporal cell action. In the positive-tail family
of recording (133), \(R_B\) grows as \(\epsilon^{-2}\). If reservoir
costs additionally obey \(M-R_B\le C\) uniformly, then
\(|P-\Phi_B|\le |P|C/R_B\) and the bridge becomes asymptotically
dominant. That extra bound must be checked; exact vacuum is not such a
check.

\subsection{3. Continuity does not close the full field
equation}\label{continuity-does-not-close-the-full-field-equation}

Equations (8) are exact for the restricted model. They do not imply that
its fixed density family is invariant under (4)'s Hamilton--Jacobi
equation. For a Gaussian base \(\rho_A=N(0,A)\) and initial \(w=P=0\),
the full equation gives
\[\partial_tS_q=\frac{\kappa^2 q}{mA^2}.\tag{11}\] The family (7)
instead permits only a multiple of \(G/\rho_A\) at that initial time. It
vanishes outside a compact interval, whereas (11) does not. This is an
immediate nonzero transverse term for \(\kappa>0\). It cannot be removed
by a spatially constant action gauge.

Adding Gaussian width supplies the missing global quadratic phase at
\(w=0\). The next calculation tests whether that repairs closure for
nonzero population transfer.

\textbf{Proposition 3 (a fixed compact transfer shape cannot close).}
Fix \(m,\kappa,A>0\) and a nonzero \(h\in C_c^\infty(\mathbb R)\) with
\(\int h=0\). The density family \(\rho_{A,w}=\rho_A+wh\), allowing
width variation and sufficiently small positive and negative \(w\),
cannot be locally invariant under (4) while including
zero-phase-gradient states for every such \(w\), satisfying continuity
with vanishing exterior flux, and having twice differentiable parameter
trajectories. The same obstruction persists with Gaussian mean motion
allowed. The claim concerns an open family of preparations, not all
finite-dimensional models.

\emph{Proof.} Set \(r=\rho_A\) and \(G=-\int_{-\infty}^q h\), so \(G\)
is smooth and compactly supported. The integrated density tangents for
width and transfer are \(G_A=qr/(2A)\) and \(G\). At an initially zero
phase gradient, continuity first gives \(\dot A=\dot w=0\).
Differentiating it and using (4) requires \[\begin{aligned}
\rho_wF_w&=c_A(w)G_A+c_w(w)G,\\
F_w&=\frac{2\kappa^2}{m}\partial_q
      \frac{(\sqrt{\rho_w})_{qq}}{\sqrt{\rho_w}},
\qquad c_A=m\ddot A,\quad c_w=m\ddot w.
\end{aligned}\tag{12}\] Outside the compact supports, \(\rho_w=r\) and
\(F_w=\kappa^2q/(mA^2)\); hence \(c_A(w)=2\kappa^2/(mA)\) for every
small \(w\). At \(w=0\) this already matches the full expression, giving
\(c_w(0)=0\). Evaluation at any point with \(G\ne0\) makes \(c_w(w)\)
smooth, since the left side is smooth while \(\rho_w>0\).

Put \(u=h/r\). Differentiating the square-root ratio gives
\(\delta[(\sqrt\rho)_{qq}/\sqrt\rho]
=(u''-qu'/A)/2\). Retaining the variation of the density multiplier in
(12), the necessary linearized identity is
\[u'''-\frac qA u''-\frac{u'}A+\frac q{A^2}u=cg,
\qquad g=G/r,\qquad c=mc_w'(0)/\kappa^2.\tag{13}\] Since \(G'=-h\), one
has \(u=qg/A-g'\). Substitution yields \[\begin{aligned}
-g''''+\frac{2q}{A}g'''
 +\left(\frac4A-\frac{q^2}{A^2}\right)g''
 -\frac{4q}{A^2}g'\\
 +\left(\frac{q^2}{A^3}-\frac1{A^2}\right)g&=cg.
\end{aligned}\tag{14}\] The function \(g\) is smooth and compactly
supported because \(r>0\). All four initial data vanish at a point
outside its support. Uniqueness for this regular fourth-order linear ODE
forces \(g\equiv0\), hence \(h=-G'\equiv0\), a contradiction. Allowing
the Gaussian mean adds a term proportional to \(r\) in (12); tail
matching forces its coefficient to zero and leaves the same
contradiction. \(\square\)

The zero-phase-gradient and continuity assumptions are essential to the
initial tangency test. The theorem permits isolated states, shapes
depending on parameters or time, and other finite-dimensional families.
It specifically rules out treating a fixed compact transfer bump plus
Gaussian width as an exact free Fisher evolution. Generated shape
information cannot be hidden inside that projection.

\subsection{4. Consequence for STATE}\label{consequence-for-state}

The canonical Fisher structure follows from a specified restriction of
ordinary classical sheet action; physical preservation of that
restriction remains open. Population memory must carry its full
transport-conjugate action, including reservoirs, before taking a
weak-bridge limit. The continuity-compatible two-coordinate family fails
the full field equation, and adding Gaussian width does not repair its
fixed compact transfer shape. Next lemma: a paired density/phase shape
mode on the expanding Gaussian, first with exact linearized field
evolution and then a positive nonlinear completion. Its tail variation
must be retained; linear mode closure alone is not a physical
population-history model. Reject a projected residual counted as exact
closure, bridge phase used as population momentum without reservoir
costs, or an unexcluded zero branch presented as positive action
necessity.
\end{document}
